%%=====================================================================
%%  RT-Comops Backend — Présentation de Soutenance (YowAuth)
%%  Format : LaTeX Beamer (Compilable sur Overleaf)
%%  Auteur : Tchoutzine Tchetnkou albino .C
%%=====================================================================
\documentclass[10pt,aspectratio=169]{beamer}

% Encodage & Langue
\usepackage[utf8]{inputenc}
\usepackage[T1]{fontenc}
\usepackage[french]{babel}
\usepackage{booktabs}
\usepackage{listings}
\usepackage{xcolor}

% Thème Beamer professionnel et épuré
\usetheme{Madrid}
\usecolortheme{default}

% Personnalisation de la charte de couleurs (Navy / Bleu ENSPY)
\definecolor{enspynavy}{HTML}{1F3A5F}
\definecolor{enspygold}{HTML}{D4AF37}
\definecolor{rtlight}{HTML}{F3F4F6}
\definecolor{rtcode}{HTML}{0B1021}

\setbeamercolor{palette primary}{bg=enspynavy,fg=white}
\setbeamercolor{palette secondary}{bg=enspynavy!85,fg=white}
\setbeamercolor{palette tertiary}{bg=enspynavy!70,fg=white}
\setbeamercolor{title}{bg=enspynavy,fg=white}
\setbeamercolor{structure}{fg=enspynavy}
\setbeamercolor{block title}{bg=enspynavy,fg=white}
\setbeamercolor{block body}{bg=rtlight,fg=black}

% Suppression des boutons de navigation Beamer en bas
\setbeamertemplate{navigation symbols}{}

% Métadonnées de la présentation
\title[Présentation YowAuth (auth-core)]{RT-Comops Backend -- Module YowAuth (auth-core)}
\subtitle{Fournisseur d'Identité (IdP) Central \& SSO}
\author[Tchoutzine Tchetnkou]{Tchoutzine Tchetnkou albino .C (Matricule : 22P387)}
\institute[ENSPY]{
  Département de Génie Informatique\\
  \textbf{École Nationale Supérieure Polytechnique de Yaoundé (ENSPY)}\\
  \vspace{0.2cm}
  \scalebox{0.85}{UE : Réseaux Mobiles et Intelligents}
}
\date{\today}

% Métadonnées additionnelles pour la page de garde
\titlegraphic{\includegraphics[height=1.2cm]{logo_enspy.png}}

\lstdefinestyle{shell}{
  backgroundcolor=\color{rtcode}, basicstyle=\ttfamily\tiny\color{white},
  breaklines=true, columns=fullflexible, frame=none,
  showstringspaces=false, aboveskip=4pt, belowskip=4pt}

\begin{document}

%---------------------------------------------------------------------
% Page de titre
%---------------------------------------------------------------------
\begin{frame}
  \titlepage
  \centering
  \scalebox{0.8}{\textcolor{enspynavy}{\textbf{Superviseur :}} Pr. Eng Thomas Djotio \quad \textbf{Manager/Chef de projet :}} \scalebox{0.8}{M. Tsafack Fotso Savio}
\end{frame}

%---------------------------------------------------------------------
% Slide 1 : Introduction et Objet du Module
%---------------------------------------------------------------------
\begin{frame}{Introduction et Rôle de YowAuth}
  \begin{block}{Qu'est-ce que YowAuth (auth-core) ?}
    C'est le module de gestion d'identité (IdP) et de Single Sign-On (SSO) de la plateforme RT-Comops Backend.
  \end{block}
  \vspace{0.2cm}
  \begin{columns}[T]
    \begin{column}{0.48\textwidth}
      \textbf{Séparation des responsabilités}
      \begin{itemize}
        \item \textbf{Identity Profile} (\code{actor-core}) : Gère les informations physiques et l'état civil de l'humain.
        \item \textbf{Authentication account} (\code{auth-core}) : Gère les identifiants techniques et les secrets d'accès.
      end{itemize}
    \end{column}
    \begin{column}{0.48\textwidth}
      \textbf{Périmètre Fonctionnel}
      \begin{itemize}
        \item Inscription publique \& Connexion locale.
        \item Résolution de contextes multi-tenant.
        \item Double authentification (MFA).
        \item Émission de tokens d'accès JWT RS256.
      \end{itemize}
    \end{column}
  \end{columns}
\end{frame}

%---------------------------------------------------------------------
% Slide 2 : Architecture Hexagonale
%---------------------------------------------------------------------
\begin{frame}{Architecture Hexagonale du Module}
  \begin{columns}[T]
    \begin{column}{0.5\textwidth}
      \textbf{Structure Interne (Hexagone)}
      \begin{itemize}
        \item \textbf{Couche de Domaine} : Contient l'agrégat immuable \code{UserAccount} et ses invariants d'intégrité.
        \item \textbf{Couche d'Application} : Coordonne les cas d'usage via les ports d'entrée et de sortie.
        \item \textbf{Couche d'Adaptateurs} : WebFlux (adaptateurs d'entrée) et R2DBC réactif (adaptateurs de sortie).
      end{itemize}
    \end{column}
    \begin{column}{0.46\textwidth}
      \begin{block}{Invariants de Domaine}
        \begin{itemize}
          \item Identifiants uniques et normalisés.
          \item Mot de passe haché avec BCrypt.
          \item Interdiction d'émettre un JWT actif si le challenge MFA OTP est en attente.
          \item Blocage automatique des comptes marqués comme suspendus ou désactivés.
        \end{itemize}
      \end{block}
    \end{column}
  \end{columns}
\end{frame}

%---------------------------------------------------------------------
% Slide 3 : Surface d'API exposée
%---------------------------------------------------------------------
\begin{frame}{Surface d'API et Sécurité}
  \textbf{Principales routes exposées (secured by backend API-Key)}
  \vspace{0.2cm}
  \begin{table}
    \small
    \begin{tabular}{lll}
      \toprule
      \textbf{Méthode} & \textbf{Endpoint} & \textbf{Rôle technique} \\
      \midrule
      \code{POST} & \code{/api/auth/identify} & Détection du type de compte (2 étapes) \\
      \code{POST} & \code{/api/auth/login} & Authentification mot de passe / Défi MFA \\
      \code{POST} & \code{/api/auth/login/mfa/confirm} & Validation du code OTP \\
      \code{POST} & \code{/api/auth/discover-contexts} & Lister les tenants associés \\
      \code{POST} & \code{/api/auth/select-context} & Émission du JWT final du contexte choisi \\
      \code{PUT} & \code{/api/users/me/plan} & Changement de forfait (Self-service) \\
      \bottomrule
    \end{tabular}
  \end{table}
\end{frame}

%---------------------------------------------------------------------
% Slide 4 : Sécurité Cryptographique et OIDC
%---------------------------------------------------------------------
\begin{frame}{Sécurité Cryptographique et OIDC}
  \begin{columns}[T]
    \begin{column}{0.48\textwidth}
      \begin{block}{Signature Asymétrique RS256}
        \begin{itemize}
          \item Les JWT sont signés avec la clé privée du Kernel.
          \item Validation décentralisée (hors-ligne) par les microservices via la clé publique.
          \item Publication de la clé publique au format standard JWKS (\code{/.well-known/jwks.json}).
        \end{itemize}
      \end{block}
    \end{column}
    \begin{column}{0.48\textwidth}
      \begin{block}{OIDC et Token Exchange (RFC 8693)}
        \begin{itemize}
          \item SSO basé sur l'échange de jetons de session.
          \item Route \code{/oauth2/token} : convertit le jeton SSO en jeton de service contextualisé (ex: \code{SALES}).
          \item Route \code{/oauth2/userinfo} : résout l'identité de l'utilisateur final.
        \end{itemize}
      \end{block}
    \end{column}
  \end{columns}
\end{frame}

%---------------------------------------------------------------------
% Slide 5 : Connexion en 2 étapes & MFA
%---------------------------------------------------------------------
\begin{frame}{Expérience utilisateur : Connexion en 2 étapes \& MFA}
  \begin{columns}[T]
    \begin{column}{0.55\textwidth}
      \textbf{Cinématique ergonomique et sécurisée}
      \begin{enumerate}
        \item \textbf{Détection de l'identifiant} : L'application interroge \code{identify} pour savoir s'il s'agit d'une inscription ou d'une connexion.
        \item \textbf{Vérification du mot de passe} : Saisie du secret. Si le compte a le MFA activé, le serveur retourne un code \code{202 Accepted} contenant un \code{mfaToken}.
        \item \textbf{Défi OTP} : L'utilisateur reçoit un OTP par e-mail et le saisit. La session est validée via \code{/login/mfa/confirm}.
      \end{enumerate}
    \end{column}
    \begin{column}{0.4\textwidth}
      \begin{block}{MFA à la volée (Self-service)}
        L'utilisateur connecté peut lui-même :
        \begin{itemize}
          \item Activer le MFA (\code{mfa/enable}).
          \item Valider son activation par OTP (\code{mfa/confirm}).
          \item Désactiver son MFA (\code{mfa/disable}).
        \end{itemize}
      \end{block}
    \end{column}
  \end{columns}
\end{frame}

%---------------------------------------------------------------------
% Slide 6 : Démonstration et Sandbox (POC)
%---------------------------------------------------------------------
\begin{frame}{Démonstration du POC (Next.js)}
  \begin{block}{Validation de l'implémentation en local et en ligne}
    Une application Next.js a été connectée au Kernel YowAuth. Elle sert de bac à sable interactif.
  \end{block}
  \vspace{0.2cm}
  \begin{itemize}
    \item \textbf{Double Colonne responsive} : Grille moderne séparant les applications métiers actives du centre de profil / sécurité.
    \item \textbf{Vitrine SSO Interactive} : Un terminal intégré trace en temps réel les appels d'échange de jeton (\code{token-exchange}) et d'introspection d'identité (\code{userinfo}) pour illustrer la mécanique sous-jacente du SSO.
    \item \textbf{Images représentatives} : Remplacement complet des emojis par des visuels professionnels pour illustrer le catalogue de services.
  \end{itemize}
\end{frame}

%---------------------------------------------------------------------
% Slide 7 : Conclusion et Perspectives
%---------------------------------------------------------------------
\begin{frame}{Conclusion et Perspectives}
  \begin{block}{Objectifs Remplis}
    \begin{itemize}
      \item Architecture hexagonale réactive implémentée de bout en bout (Spring Boot WebFlux).
      \item Cycle IAM complet opérationnel (inscription, connexion multi-contexte, MFA).
      \item Standardisation sur OIDC / Token Exchange et cryptographie RS256/JWKS.
    \end{itemize}
  \end{block}
  \vspace{0.2cm}
  \textbf{Perspectives de développement}
  \begin{itemize}
    \item Support des clés physiques FIDO2 / WebAuthn pour le MFA sans mot de passe.
    \item Liaison directe avec un serveur SMTP en production (fin du mode \code{PREVIEW_ONLY}).
    \item Back-office de gestion et de suspension administrative des comptes utilisateurs.
  \end{itemize}
\end{frame}

%---------------------------------------------------------------------
% Page finale
%---------------------------------------------------------------------
\begin{frame}
  \centering
  \Huge \textbf{\textcolor{enspynavy}{Merci pour votre attention !}}
  \vspace{0.8cm}
  
  \Large Questions \& Réponses
  \vspace{0.8cm}
  
  \scalebox{0.85}{\textcolor{rtgray}{École Nationale Supérieure Polytechnique de Yaoundé (ENSPY)}}
\end{frame}

\end{document}
